\documentclass[11pt,a4paper]{article}
\usepackage[top=2.3cm,bottom=2.3cm,left=2.2cm,right=2.2cm]{geometry}

\usepackage{amsmath,amssymb}
\usepackage{fontspec}
\setmainfont{Liberation Serif}
\setsansfont{Liberation Sans}
\setmonofont{DejaVu Sans Mono}[Scale=0.84]

\usepackage{xcolor}
\definecolor{brandblue}{RGB}{20, 55, 105}
\definecolor{accentblue}{RGB}{35, 95, 165}
\definecolor{codebg}{RGB}{248, 249, 251}
\definecolor{codeframe}{RGB}{215, 222, 230}
\definecolor{javakeyword}{RGB}{0, 45, 170}
\definecolor{javastring}{RGB}{5, 120, 20}
\definecolor{javacomment}{RGB}{120, 120, 120}
\definecolor{javanumber}{RGB}{20, 75, 220}

\definecolor{noteborder}{RGB}{30, 110, 65}
\definecolor{notebg}{RGB}{242, 249, 245}
\definecolor{warnborder}{RGB}{185, 75, 10}
\definecolor{warnbg}{RGB}{254, 248, 240}
\definecolor{tipborder}{RGB}{30, 95, 160}
\definecolor{tipbg}{RGB}{242, 247, 254}

\usepackage{fancyhdr}
\usepackage{lastpage}
\pagestyle{fancy}
\fancyhf{}
\lhead{\parbox[t]{0.58\textwidth}{\raggedright\small\sffamily\color{brandblue}\textbf{T09 --- Array, Matrici e Tipi Enumerativi (enum)}}}
\rhead{\small\sffamily\color{gray}Programmazione II -- UniVR (2025/2026)}
\cfoot{\small\sffamily Pagina \thepage\ di \pageref{LastPage}}
\renewcommand{\headrulewidth}{0.5pt}
\renewcommand{\headrule}{\hbox to\headwidth{\color{brandblue!70!black}\leaders\hrule height \headrulewidth\hfill}}

\usepackage{titlesec}
\titleformat{\section}{\Large\bfseries\sffamily\color{brandblue}}{\thesection}{0.8em}{}[\titlerule]
\titleformat{\subsection}{\large\bfseries\sffamily\color{accentblue}}{\thesubsection}{0.8em}{}
\titleformat{\subsubsection}{\normalsize\bfseries\sffamily\color{brandblue!85!black}}{\thesubsubsection}{0.8em}{}

\usepackage{needspace}
\usepackage{fancyvrb}
\DefineVerbatimEnvironment{asciibox}{Verbatim}{
  fontsize=\fontsize{7.5pt}{9.2pt}\selectfont,
  frame=single,
  rulecolor=\color{codeframe},
  fillcolor=\color{codebg},
  framesep=2.5mm
}
\usepackage{listings}
\lstset{
  backgroundcolor=\color{codebg},
  frame=single,
  rulecolor=\color{codeframe},
  basicstyle=\ttfamily\footnotesize,
  keywordstyle=\color{javakeyword}\bfseries,
  stringstyle=\color{javastring},
  commentstyle=\color{javacomment}\itshape,
  numberstyle=\tiny\color{gray},
  numbers=left,
  numbersep=7pt,
  stepnumber=1,
  breaklines=true,
  breakatwhitespace=true,
  showstringspaces=false,
  tabsize=4,
  xleftmargin=12pt,
  framexleftmargin=12pt
}

\newcommand{\passthrough}[1]{#1}
\providecommand{\tightlist}{\setlength{\itemsep}{0pt}\setlength{\parskip}{0pt}}

\usepackage{tcolorbox}
\tcbuselibrary{skins,breakable}
\newtcolorbox{studynote}[1][Nota]{
  enhanced,
  breakable,
  colback=notebg,
  colframe=noteborder,
  fonttitle=\bfseries\sffamily\small,
  coltitle=white,
  title={#1},
  arc=2mm,
  boxrule=0.8pt,
  left=9pt, right=9pt, top=7pt, bottom=7pt
}

\newtcolorbox{studywarning}[1][Attenzione]{
  enhanced,
  breakable,
  colback=warnbg,
  colframe=warnborder,
  fonttitle=\bfseries\sffamily\small,
  coltitle=white,
  title={#1},
  arc=2mm,
  boxrule=0.8pt,
  left=9pt, right=9pt, top=7pt, bottom=7pt
}

\newtcolorbox{studytip}[1][Suggerimento]{
  enhanced,
  breakable,
  colback=tipbg,
  colframe=tipborder,
  fonttitle=\bfseries\sffamily\small,
  coltitle=white,
  title={#1},
  arc=2mm,
  boxrule=0.8pt,
  left=9pt, right=9pt, top=7pt, bottom=7pt
}

\usepackage{booktabs}
\usepackage{longtable}
\usepackage{array}
\usepackage{calc}

\usepackage{hyperref}
\hypersetup{
  colorlinks=true,
  linkcolor=accentblue,
  urlcolor=accentblue,
  citecolor=brandblue,
  pdfborder={0 0 0}
}

\title{\Huge\bfseries\sffamily\color{brandblue} T09 --- Array, Matrici e Tipi Enumerativi (enum) \\[0.3em] \large\sffamily Array Mono e Multidimensionali, Enum Tipizzati, Immutabilità e Singleton in SimpleDate v4}
\author{\textbf{Docente:} Prof. Michele Pasqua \quad---\quad \textbf{Appunti Revisione:} Wally}
\date{A.A. 2025/2026 -- Universit\`a degli Studi di Verona}

\begin{document}
\maketitle
\thispagestyle{fancy}
\vspace{-1.2em}
\noindent\textcolor{brandblue}{\rule{\textwidth}{0.8pt}}
\vspace{1.2em}

In questo blocco approfondiamo la gestione a basso livello degli array
nello Heap, i tipi enumerati (\passthrough{\lstinline!enum!})
(\href{file:///home/wally/Documenti/GitHub/Programmazione-II/25-26/Teoria-e-Esercitazioni/Slides-20260902/T09\%20-\%20Arrays\%20and\%20Enumerative\%20Types.pdf}{T09
- Arrays and Enumerative Types.pdf}) e analizziamo il codice di
\href{file:///home/wally/Documenti/GitHub/Programmazione-II/25-26/Teoria-e-Esercitazioni/Codice-20260902/Lezione09/SimpleDate/src}{\passthrough{\lstinline!Lezione09!}},
dove \passthrough{\lstinline!Date!} diventa una \textbf{classe
immutabile} (campi \passthrough{\lstinline!final!} e rimozione dei
setter) e compare per la prima volta un design pattern fondamentale: il
\textbf{Singleton Pattern} implementato in
\href{file:///home/wally/Documenti/GitHub/Programmazione-II/25-26/Teoria-e-Esercitazioni/Codice-20260902/Lezione09/SimpleDate/src/BirthDay.java}{\passthrough{\lstinline!BirthDay.java!}}.

\begin{center}\rule{0.5\linewidth}{0.5pt}\end{center}

\subsubsection{1. Teoria: Concetti Teorici dalle Slide (Slide
T09)}\label{teoria-concetti-teorici-dalle-slide-slide-t09}

\paragraph{A. Gli Array in Java (Slide
1--13)}\label{a.-gli-array-in-java-slide-113}

\begin{itemize}
\item
  \textbf{Natura degli Array}:

  \begin{itemize}
  \tightlist
  \item
    Un array è una sequenza contigua e ordinata di variabili omogenee
    (dello stesso tipo), accessibili tramite indice numerico intero che
    parte da \passthrough{\lstinline!0!}.
  \item
    In Java gli array \textbf{non sono tipi primitivi}: sono
    \textbf{oggetti speciali allocati nello Heap}.
  \item
    Possono memorizzare valori primitivi oppure
    \textbf{riferimenti/puntatori a oggetti}, mai oggetti per valore
    diretto.
  \item
    La dimensione di un array può essere stabilita dinamicamente a
    runtime all'atto dell'allocazione
    (\passthrough{\lstinline!new int[size]!}), ma \textbf{è immutabile
    una volta creata}.
  \end{itemize}
\item
  \textbf{Dichiarazione vs Creazione}:

  \begin{itemize}
  \tightlist
  \item
    Sintassi raccomandata: \passthrough{\lstinline!int[] arr;!} (la
    notazione \passthrough{\lstinline!int arr[];!} è ammessa per
    retrocompatibilità col C, ma sconsigliata).
  \item
    La dichiarazione \passthrough{\lstinline!int[] arr;!} alloca solo
    una variabile reference nello Stack (inizializzata a
    \passthrough{\lstinline!null!}), \textbf{senza riservare spazio per
    gli elementi nello Heap}.
  \item
    Allocazione:

    \begin{itemize}
    \tightlist
    \item
      Tramite \passthrough{\lstinline!new!}:
      \passthrough{\lstinline!float[] arr = new float[10];!} (gli
      elementi assumono il valore di default del tipo, es.
      \passthrough{\lstinline!0.0f!}).
    \item
      Tramite inizializzazione statica:
      \passthrough{\lstinline!int[] primes = \{2, 3, 5, 7\};!} (la
      dimensione viene inferita automaticamente dal compilatore).
    \end{itemize}
  \end{itemize}
\item
  \textbf{Il Pseudo-Campo \passthrough{\lstinline!.length!}}:

  \begin{itemize}
  \tightlist
  \item
    Ogni array possiede la proprietà
    \passthrough{\lstinline!public final int length!}.
  \item
    Attenzione: \textbf{Distinzione d'esame}: per gli array si scrive
    \passthrough{\lstinline!arr.length!} (\textbf{senza parentesi}, è
    una variabile \passthrough{\lstinline!final!}), mentre per le
    stringhe si invoca il metodo \passthrough{\lstinline!str.length()!}
    (\textbf{con le parentesi}).
  \item
    Se si tenta di leggere \passthrough{\lstinline!arr.length!} su un
    array nullo (\passthrough{\lstinline!arr = null;!}), la JVM lancia
    una \passthrough{\lstinline!NullPointerException!}.
  \end{itemize}
\item
  \textbf{Confronto e Stampa degli Array (I Trabocchetti Classici)}:

  \begin{enumerate}
  \def\labelenumi{\arabic{enumi}.}
  \tightlist
  \item
    \passthrough{\lstinline!arr1 == arr2!}: confronta solo gli indirizzi
    nello Heap. Se due array contengono gli stessi identici elementi ma
    risiedono in aree diverse, restituisce
    \passthrough{\lstinline!false!}.
  \item
    \passthrough{\lstinline!arr1.equals(arr2)!}: \textbf{NON funziona}.
    Gli array non sovrascrivono il metodo
    \passthrough{\lstinline!.equals()!} di
    \passthrough{\lstinline!Object!}, quindi
    \passthrough{\lstinline!equals!} si comporta esattamente come
    \passthrough{\lstinline!==!}.
  \item
    \textbf{Soluzione standard}: per confrontare il contenuto si usa il
    metodo statico
    \textbf{\passthrough{\lstinline!java.util.Arrays.equals(arr1, arr2)!}}.
  \item
    \passthrough{\lstinline!arr.toString()!}: stampa la firma bytecode
    della reference (es. \passthrough{\lstinline![I@6ce253f1!}, dove
    \passthrough{\lstinline![!} indica un array e
    \passthrough{\lstinline!I!} indica il tipo
    \passthrough{\lstinline!int!}).
  \item
    Per visualizzare gli elementi a video in formato leggibile:
    \textbf{\passthrough{\lstinline!java.util.Arrays.toString(arr)!}}.
  \end{enumerate}
\item
  \textbf{Enhanced For Loop (For-Each, Java 5+)}:

\begin{lstlisting}[language=Java]
for (String arg : args) {
    System.out.println(arg);
}
\end{lstlisting}

  Sostituisce il for contatore eliminando il rischio di errori di
  \emph{off-by-one} e \emph{ArrayIndexOutOfBoundsException}.
\item
  \textbf{Array Multidimensionali e Matrici Frastagliate (\emph{Jagged
  Arrays})}:

  \begin{itemize}
  \tightlist
  \item
    In Java non esistono matrici bidimensionali a blocco contiguo in
    stile Fortran/C: \textbf{un array multidimensionale è un array di
    riferimenti ad altri array}.
  \item
    \emph{Conseguenze pratiche}:

    \begin{enumerate}
    \def\labelenumi{\arabic{enumi}.}
    \item
      \textbf{Scambio di righe in tempo \(O(1)\)}: per invertire la
      prima e l'ultima riga di una matrice non serve copiare gli
      elementi uno per uno, basta scambiare i loro puntatori:

\begin{lstlisting}[language=Java]
int[] temp = matrix[0];
matrix[0] = matrix[matrix.length - 1];
matrix[matrix.length - 1] = temp;
\end{lstlisting}
    \item
      \textbf{Righe a lunghezza eterogenea (\emph{Jagged/Pyramid
      Arrays})}:

\begin{lstlisting}[language=Java]
int[][] pyramid = new int[3][];
pyramid[0] = new int[1]; // Riga 0 ha 1 elemento
pyramid[1] = new int[2]; // Riga 1 ha 2 elementi
pyramid[2] = new int[3]; // Riga 2 ha 3 elementi
\end{lstlisting}
    \item
      Per stampare matrici annidate,
      \passthrough{\lstinline!Arrays.toString(matrix)!} stampa gli
      indirizzi delle singole righe. Si deve iterare sulle righe o usare
      \textbf{\passthrough{\lstinline!java.util.Arrays.deepToString(matrix)!}}.
    \end{enumerate}
  \end{itemize}
\end{itemize}

\begin{center}\rule{0.5\linewidth}{0.5pt}\end{center}

\paragraph{\texorpdfstring{B. Tipi Enumerati (\texttt{enum}) (Slide
14--19)}{B. Tipi Enumerati (enum) (Slide 14--19)}}\label{b.-tipi-enumerati-enum-slide-1419}

\begin{itemize}
\tightlist
\item
  Introdotti in Java 5 con la keyword \passthrough{\lstinline!enum!} per
  gestire insiemi chiusi e prefissati di valori costanti (es. giorni
  della settimana, punti cardinali, lingue).
\item
  \textbf{Regole e Caratteristiche}:

  \begin{itemize}
  \tightlist
  \item
    Possono essere dichiarati in un file \passthrough{\lstinline!.java!}
    dedicato (come classe autonoma) o all'interno di una classe come
    membro statico, ma \textbf{mai all'interno di un metodo}.
  \item
    Non sono semplici interi o stringhe: sono classi speciali che
    estendono implicitamente \passthrough{\lstinline!java.lang.Enum!}.
    Il loro costruttore è privato e non può essere invocato con
    \passthrough{\lstinline!new!}.
  \item
    \textbf{Metodi Predefiniti Fondamentali}:

    \begin{itemize}
    \tightlist
    \item
      \passthrough{\lstinline!e.name()!}: restituisce il nome letterale
      della costante come \passthrough{\lstinline!String!} (es.
      \passthrough{\lstinline!"WEST"!}).
    \item
      \passthrough{\lstinline!e.ordinal()!}: restituisce la posizione
      ordinale intera partendo da 0 (es. \passthrough{\lstinline!3!}).
    \item
      \passthrough{\lstinline!EnumClass.values()!}: restituisce un array
      contenente tutte le costanti dichiarate nell'enum, comodo per
      cicli for-each.
    \end{itemize}
  \end{itemize}
\item
  \textbf{Confronto tra Enum: Perché usare \passthrough{\lstinline!==!}
  invece di \passthrough{\lstinline!.equals()!}}:

  \begin{itemize}
  \tightlist
  \item
    La JVM garantisce che in memoria esista \textbf{una e una sola
    istanza} per ciascuna costante di un enum (\emph{garanzia di unicità
    singleton}).
  \item
    Pertanto, il confronto con \textbf{\passthrough{\lstinline!==!}} è
    perfettamente sicuro ed è \textbf{più robusto di
    \passthrough{\lstinline!.equals()!}}: se una variabile vale
    \passthrough{\lstinline!null!}, l'espressione
    \passthrough{\lstinline!move == Direction.WEST!} restituisce
    \passthrough{\lstinline!false!} senza errori, mentre
    \passthrough{\lstinline!move.equals(...)!} causerebbe un crash con
    \passthrough{\lstinline!NullPointerException!}!
  \end{itemize}
\item
  \textbf{Uso di Enum negli \passthrough{\lstinline!switch!}}:

  \begin{itemize}
  \tightlist
  \item
    All'interno dei blocchi \passthrough{\lstinline!case!}, \textbf{non
    si qualifica} il tipo: si scrive \passthrough{\lstinline!case IT:!}
    e non \passthrough{\lstinline!case Language.IT:!}.
  \end{itemize}
\end{itemize}

\begin{center}\rule{0.5\linewidth}{0.5pt}\end{center}

\subsubsection{\texorpdfstring{2. Codice: Evoluzione del Codice:
\texttt{SimpleDate} (Lezione
09)}{2. Codice: Evoluzione del Codice: SimpleDate (Lezione 09)}}\label{codice-evoluzione-del-codice-simpledate-lezione-09}

In
\href{file:///home/wally/Documenti/GitHub/Programmazione-II/25-26/Teoria-e-Esercitazioni/Codice-20260902/Lezione09/SimpleDate/src}{\passthrough{\lstinline!Lezione09!}},
il docente introduce due cambiamenti strutturali fondamentali:

\paragraph{\texorpdfstring{1. Immutabilità di
\href{file:///home/wally/Documenti/GitHub/Programmazione-II/25-26/Teoria-e-Esercitazioni/Codice-20260902/Lezione09/SimpleDate/src/Date.java}{\texttt{Date.java}}}{1. Immutabilità di Date.java}}\label{immutabilituxe0-di-date.java}

\begin{lstlisting}
 public class Date {
-    private int day;
-    private int month;
-    private int year;
+    private final int day;
+    private final int month;
+    private final int year;

     // Costruttori...

-    public void setDay(int day) { ... }
-    public void setMonth(int month) { ... }
-    public void setYear(int year) { ... }
\end{lstlisting}

\begin{itemize}
\tightlist
\item
  \textbf{Campi \passthrough{\lstinline!final!}}: Una volta
  inizializzati nel costruttore, \passthrough{\lstinline!day!},
  \passthrough{\lstinline!month!} e \passthrough{\lstinline!year!} non
  possono più essere riassegnati.
\item
  \textbf{Rimozione dei Setter}: I metodi mutatori
  \passthrough{\lstinline!setDay!}, \passthrough{\lstinline!setMonth!} e
  \passthrough{\lstinline!setYear!} vengono eliminati.
\item
  \textbf{Pattern Immutabile}: Un oggetto \passthrough{\lstinline!Date!}
  è ora a prova di manomissione. Se un programma vuole calcolare ``il
  giorno successivo'', non può mutare l'oggetto esistente, ma
  \textbf{deve instanziare un nuovo oggetto
  \passthrough{\lstinline!Date!}} con le coordinate aggiornate
  (esattamente come avviene per le \passthrough{\lstinline!String!}).
\end{itemize}

\begin{center}\rule{0.5\linewidth}{0.5pt}\end{center}

\paragraph{\texorpdfstring{2. Il Design Pattern Singleton:
\href{file:///home/wally/Documenti/GitHub/Programmazione-II/25-26/Teoria-e-Esercitazioni/Codice-20260902/Lezione09/SimpleDate/src/BirthDay.java}{\texttt{BirthDay.java}}}{2. Il Design Pattern Singleton: BirthDay.java}}\label{il-design-pattern-singleton-birthday.java}

\begin{lstlisting}[language=Java]
public class BirthDay {
    private static Date date = new Date(1, 1, 1970);
    private static BirthDay instance; // Riferimento all'unica istanza

    // 1. Costruttore privato: impedisce istanziazioni esterne con 'new BirthDay()'
    private BirthDay() { }

    // 2. Metodo statico di fabbrica con Lazy Initialization
    public static BirthDay getInstance() {
        if (instance == null) 
            instance = new BirthDay();
        return instance;
    }

    public Date getDate() { 
        return date; 
    }
}
\end{lstlisting}

\paragraph{Analisi: Analisi Architetturale del
Singleton:}\label{analisi-analisi-architetturale-del-singleton}

\begin{itemize}
\tightlist
\item
  \textbf{Scopo del Pattern}: Garantire che per tutta la durata
  dell'applicazione esista \textbf{una e una sola istanza} di una classe
  nello Heap e fornire un punto di accesso globale ad essa.
\item
  \textbf{I Tre Pilastri del Singleton in Java}:

  \begin{enumerate}
  \def\labelenumi{\arabic{enumi}.}
  \tightlist
  \item
    \textbf{Costruttore \passthrough{\lstinline!private!}}: blocca
    l'accesso a \passthrough{\lstinline!new BirthDay()!} da qualsiasi
    altra classe (tentare di chiamarlo genera un errore a compile-time).
  \item
    \textbf{Campo statico privato (\passthrough{\lstinline!instance!})}:
    contiene l'unico riferimento all'oggetto.
  \item
    \textbf{Metodo factory pubblico statico
    (\passthrough{\lstinline!getInstance()!})}: implementa la \emph{Lazy
    Initialization} (crea l'oggetto solo alla prima richiesta,
    riusandolo per tutte le successive).
  \end{enumerate}
\end{itemize}

\begin{center}\rule{0.5\linewidth}{0.5pt}\end{center}

\paragraph{\texorpdfstring{3. Verifica in
\href{file:///home/wally/Documenti/GitHub/Programmazione-II/25-26/Teoria-e-Esercitazioni/Codice-20260902/Lezione09/SimpleDate/src/MainDate.java}{\texttt{MainDate.java}}
e il Bug dei
Getter}{3. Verifica in MainDate.java e il Bug dei Getter}}\label{verifica-in-maindate.java-e-il-bug-dei-getter}

\begin{lstlisting}[language=Java]
BirthDay bDay = BirthDay.getInstance();
// BirthDay day1 = new BirthDay(); // COMPILE-TIME ERROR!
BirthDay day1 = BirthDay.getInstance();
System.out.println("bDay ?= day1: " + (bDay == day1)); // STAMPA TRUE!

Date today = new Date(27, 10, 2025);
// today.setDay(today.getDay() + 1); // COMPILE-TIME ERROR: l'oggetto è immutabile!
Date tomorrow = new Date(today.getDay() + 1, today.getMonth(), today.getYear());
\end{lstlisting}

\begin{studywarning}[Attenzione / Trappola d'Esame] \textbf{Il Bug Persistente dei Getter e l'Effetto a
Runtime}: Anche in \passthrough{\lstinline!Lezione09!},
\passthrough{\lstinline!Date.java!} ha ancora il refuso di copia-incolla
alle righe 43--44:

\begin{lstlisting}[language=Java]
public int getMonth() { return day; } // Restituisce il campo day!
public int getYear() { return day; }  // Restituisce il campo day!
\end{lstlisting}

Quando \passthrough{\lstinline!MainDate!} crea
\passthrough{\lstinline!tomorrow!} invocando
\passthrough{\lstinline!today.getMonth()!} su una data con
\passthrough{\lstinline!day = 27!}, ottiene \passthrough{\lstinline!27!}
come mese! Di conseguenza, il costruttore riceve
\passthrough{\lstinline!(28, 27, 27)!} e
\passthrough{\lstinline!verify()!} stampa:
\passthrough{\lstinline"Illegal date!"}
\passthrough{\lstinline!tomorrow: 28/27/27!} Nel tuo codice è opportuno
correggere questi getter (\passthrough{\lstinline!return month;!} e
\passthrough{\lstinline!return year;!}).
\end{studywarning}

\begin{center}\rule{0.5\linewidth}{0.5pt}\end{center}

\subsubsection{3. Ciclo: Corrispondenza con i Tuoi Moduli
Workspace}\label{ciclo-corrispondenza-con-i-tuoi-moduli-workspace}

Nel tuo repository
\href{file:///home/wally/Documenti/GitHub/Programmazione-II/25-26/esercizi-wally-25-26}{\passthrough{\lstinline!esercizi-wally-25-26!}}:
- Nel modulo \textbf{\passthrough{\lstinline!es09!}}: hai formalizzato
questo esatto pattern creando
\href{file:///home/wally/Documenti/GitHub/Programmazione-II/25-26/esercizi-wally-25-26/es09/src/main/java/es09/SingletonPattern.java}{\passthrough{\lstinline!SingletonPattern.java!}}
con costruttore privato e metodo factory
\passthrough{\lstinline!getInstance()!}.

\begin{center}\rule{0.5\linewidth}{0.5pt}\end{center}

\begin{studynote}[Nota di Studio] Con la \textbf{Lezione 09} abbiamo consolidato la
manipolazione avanzata degli array, le proprietà degli enum,
l'immutabilità e il pattern Singleton.

Il prossimo blocco è la \textbf{Lezione 10 / Slide T10: \emph{Java
Variables and Call Stack}}: - Modello formale della memoria JVM:
\textbf{Call Stack (Stack Frames)} vs \textbf{Heap (Objects \& Class
Data)}. - Variabili d'istanza, variabili di classe
(\passthrough{\lstinline!static!}), variabili locali e parametri di
metodo. - Il passaggio dei parametri per valore (\emph{pass-by-value})
per tipi primitivi vs reference types. - Il blocco di inizializzazione
statica (\passthrough{\lstinline!static \{ ... \}!}). - \textbf{Il
grande spartiacque del corso}: il debutto del progetto
\textbf{\passthrough{\lstinline!MavenDate!}}, con l'adozione ufficiale
di Apache Maven, \passthrough{\lstinline!pom.xml!}, layout directory
standard e primo packaging!
\end{studynote}

\textbf{Dammi conferma per aprire la Lezione 10 / T10 e analizzare Call
Stack e MavenDate!}


\end{document}
